(1) 整个系统隔热，气体自由膨胀到真空室时没有推动外界活塞的边界功，因此 \(Q=W=0\)，\(\Delta U=0\)。代入所给内能函数得
\[
nC_VT_f-\frac{an^2}{V_f}
=nC_VT_i-\frac{an^2}{V_i}.
\]
于是
\[
T_f=T_i+\frac{an}{C_V}\left(\frac1{V_f}-\frac1{V_i}\right)<T_i.
\]
随体积增大，吸引势能 \(-an^2/V\) 上升，部分热运动内能转化为分子间势能，故温度下降。

(2) 对内能微分，
\[
dU=nC_VdT+\frac{an^2}{V^2}dV.
\]
对虚拟可逆状态变更使用 \(T\,dS=dU+P\,dV\)，并代入状态方程：
\[
dS=\frac{nC_V}{T}dT+
\left[\frac{an^2}{TV^2}
+\frac{nR}{V-nb}-\frac{an^2}{TV^2}\right]dV
=\frac{nC_V}{T}dT+\frac{nR}{V-nb}dV.
\]
故熵是状态函数且
\[
\Delta S=nC_V\ln\frac{T_f}{T_i}
+nR\ln\frac{V_f-nb}{V_i-nb}.
\]
此计算没有假定真实自由膨胀是可逆过程。

(3) 沿给定等内能的平衡状态族，\(dU=0\)，因此由基本关系直接得 \((\partial S/\partial V)_U=P/T\)。也可由 \(dT/dV|_U=-an/(C_VV^2)\) 代入上式独立验证：
\[
\frac{dS}{dV}\bigg|_U=
-\frac{an^2}{TV^2}+\frac{nR}{V-nb}
=\frac PT.
\]
题目明确规定在所走状态区间 \(T>0,P>0\)，故 \(dS/dV|_U>0\)，对 \(V_i<V_f\) 积分即得 \(\Delta S>0\)，与第二定律一致。这一论证只通过虚拟平衡态连接两端点，并不宣称真实膨胀过程保持空间均匀。

(4) 当 \(V_f\downarrow V_i\)，温度对最终体积的一阶斜率为
\[
\left.\frac{\partial T_f}{\partial V_f}\right|_{U,V_f=V_i}
=-\frac{an}{C_VV_i^2}<0.
\]
若先令 \(a\to0\)，体积相关吸引能消失，\(T_f\to T_i\)；进一步令 \(b\to0\)，熵变恢复 \(\Delta S=nR\ln(V_f/V_i)\)，正是理想气体向真空自由膨胀的结果。非理想气体内能含体积项，故不能不作检验地套用“绝热自由膨胀温度不变”。
